\section{Related Work}
\label{sec:related}

\paragraph{Simple baselines versus deep models.}
\citet{zhang2015character} compared bag-of-words/TF-IDF linear models with deep
convolutional networks on large text classification corpora and found that
classical models remain strong on smaller datasets, while deep models win at
millions of examples. \citet{joulin2017bag} showed that a linear model over
$n$-gram embeddings (fastText) matches deep models at a fraction of the cost.
\citet{iyyer2015deep} introduced deep averaging networks (DAN), which average
word embeddings before an MLP and rival syntax-aware models; our \texttt{emb\_mlp}
follows this design.

\paragraph{Sequence backbones.}
We compare three recurrent cells (a plain Elman RNN~\cite{elman1990finding} and the gated
LSTM~\cite{hochreiter1997long} and GRU~\cite{cho2014learning}), each unidirectional and
bidirectional~\cite{schuster1997bidirectional}, a small Transformer encoder~\cite{vaswani2017attention}
and a selective state-space model (Mamba~\cite{gu2023mamba}), all trained from
scratch with matched non-embedding parameter budgets. Tree ensembles over TF-IDF
features (Random Forest~\cite{breiman2001random}, XGBoost~\cite{chen2016xgboost})
complement the linear models.

\paragraph{Learning curves and data value.}
\citet{kolachina2012prediction} model how performance grows with training size;
we use the macro-F1 gain between adjacent sizes as a simple, non-parametric measure
of the marginal value of labels.

\paragraph{Korean review data.}
NSMC~\cite{nsmc} is the standard Korean movie-review sentiment benchmark and is
widely used in recent Korean NLP work~\cite{jang2024enhancing,han2024deep}; the
Naver Shopping and Steam review corpora~\cite{bab2min_corpus} cover product and
game reviews~\cite{park2023offensive}.
